% proofs/induction_head.tex
% Formal Proof: General Linear Position and In-Context Induction-Head Copying Error Bound

\begin{theorem}[\textbf{In-Context Induction-Head Circuit Copying Error Bound}]
\label{thm:induction_head}
Consider a 2-layer Kernel-DeltaNet architecture configured to execute an in-context associative recall task on a sequence containing $K$ associative pairs $[(x_1, y_1), \dots, (x_K, y_K)]$. Let $k_i = W_k x_i \in \mathbb{R}^{d_k}$ and $v_i = W_v y_i \in \mathbb{R}^{d_v}$ with orthogonal value representations $\mathbb{E}[v_i v_j^T] = \delta_{ij} I_{d_v}$ and unit norms $\|v_i\|_2 = 1$. Let distinct key directions satisfy maximum cosine similarity $\mu_{\max} \triangleq \max_{j \neq i} \tilde{k}_i^T \tilde{k}_j < 1$.
\begin{enumerate}[label=(\roman*)]
    \item \textbf{General Linear Position}: Under the spherically renormalized dual feature map $\hat{\phi}: \mathbb{S}^{d_k-1} \to \mathbb{S}^{D-1}$ with ORF projection $\Omega \in \mathbb{R}^{D \times d_k}$ and uniform bias $b \sim \operatorname{Uniform}[0, 2\pi]^D$, the set of lifted keys $\{\hat{\phi}(k_1), \dots, \hat{\phi}(k_K)\}$ is in general linear position in $\mathbb{R}^D$ with probability $1$, satisfying:
    \begin{equation}
    \operatorname{rank}\left( [\hat{\phi}(k_1), \dots, \hat{\phi}(k_K)] \right) = \min(K, D)
    \label{eq:glp_rank}
    \end{equation}
    \item \textbf{Multi-Token Copying Error Bound}: The expected multi-token associative copying error $\mathcal{E}_{\text{copy}} \triangleq \|y_t - v_{\text{target}}\|_2$ for query $q = k_{\text{target}}$ satisfies:
    \begin{equation}
    \mathbb{E}\left[ \mathcal{E}_{\text{copy}} \right] \le \sqrt{K-1} \left( \exp\left( -\frac{1 - \mu_{\max}}{\tau^2} \right) + \sqrt{\frac{d_k - 1}{D(d_k + 2)}} \right) = \mathcal{O}\left( \frac{K}{D} \exp\left( -\frac{2}{\tau^2} \right) \right)
    \label{eq:copy_error_bound}
    \end{equation}
    \item \textbf{Sharp Phase Transition vs. Linear DeltaNet}: For Linear DeltaNet ($D = d_k$), the copying error undergoes a catastrophic phase transition:
    \begin{equation}
    \mathcal{E}_{\text{copy}}^{\text{Linear}} = \begin{cases} \mathcal{O}\left(\sqrt{\frac{K}{d_k}}\right), & \text{if } K \le d_k \\ \Omega(1) \quad (\text{Recall Collapse}), & \text{if } K > d_k \end{cases}
    \label{eq:linear_collapse}
    \end{equation}
    whereas Kernel-DeltaNet suppresses cross-talk noise exponentially, maintaining near-zero copying error $\mathcal{E}_{\text{copy}} \le 0.05$ up to $K = D = 512$.
\end{enumerate}
\end{theorem}

\begin{proof}
\textbf{Part (i): General Linear Position via Analytic Variety Measure Zero.}
Let $\Phi \triangleq [\hat{\phi}(k_1), \dots, \hat{\phi}(k_K)] \in \mathbb{R}^{D \times K}$. A set of $K$ vectors in $\mathbb{R}^D$ fails to be in general linear position if and only if there exists a subset of $m \le \min(K, D)$ vectors that is linearly dependent. Linear dependence of $m$ vectors $\{\hat{\phi}(k_{i_1}), \dots, \hat{\phi}(k_{i_m})\}$ requires that all $m \times m$ minors of the submatrix vanish identically:
\begin{equation}
\det\left( \Phi_{\mathcal{I}, \mathcal{J}} \right) = 0, \qquad |\mathcal{I}| = |\mathcal{J}| = m
\label{eq:minor_det}
\end{equation}
The feature mapping $\hat{\phi}(k) = \frac{\cos(\Omega \tilde{k} / \tau + b)}{\|\cos(\Omega \tilde{k} / \tau + b)\|_2}$ is a composition of real-analytic trigonometric functions on distinct inputs $\tilde{k}_i \neq \tilde{k}_j$. Consequently, the determinant $\det(\Phi_{\mathcal{I}, \mathcal{J}}(\Omega, b))$ is a real-analytic function of the parameters $(\Omega, b) \in \mathbb{R}^{D \times d_k} \times [0, 2\pi]^D$. By the identity theorem for real-analytic functions, a real-analytic function on a connected domain either vanishes everywhere or has a zero locus of Lebesgue measure zero. Since the continuous RBF kernel is strictly positive definite on distinct points, the Gram matrix has non-zero determinant for some realization of $(\Omega, b)$, meaning $\det(\Phi_{\mathcal{I}, \mathcal{J}})$ does not identically vanish. Therefore, the set of parameters $(\Omega, b)$ where any $m$ vectors are linearly dependent has Lebesgue measure zero under the Haar and uniform measures. Thus, the column rank is $\min(K, D)$ with probability $1$.

\bigskip
\noindent
\textbf{Part (ii): Derivation of the Copying Error Bound.}
In Layer 1, the recurrent state accumulates the $K$ associative pairs: $S_K = \sum_{i=1}^K v_i \hat{\phi}(k_i)^T$. When queried with key $q = k_{i^*}$ (where $v_{\text{target}} \equiv v_{i^*}$):
\begin{align}
y_{\text{ret}} = S_K \hat{\phi}(k_{i^*}) &= \sum_{i=1}^K v_i \left( \hat{\phi}(k_i)^T \hat{\phi}(k_{i^*}) \right) \nonumber \\
&= v_{i^*} \|\hat{\phi}(k_{i^*})\|_2^2 + \sum_{i \neq i^*} v_i \left( \hat{\phi}(k_i)^T \hat{\phi}(k_{i^*}) \right) \nonumber \\
&= v_{\text{target}} + \sum_{i \neq i^*} v_i \alpha_i
\end{align}
where $\alpha_i \triangleq \hat{\phi}(k_i)^T \hat{\phi}(k_{i^*}) \in [-1, 1]$. The instantaneous copying error vector is $e \triangleq y_{\text{ret}} - v_{\text{target}} = \sum_{i \neq i^*} v_i \alpha_i$.
Taking the expected squared Euclidean norm:
\begin{equation}
\mathbb{E}\left[ \|e\|_2^2 \right] = \mathbb{E}\left[ \left\| \sum_{i \neq i^*} v_i \alpha_i \right\|_2^2 \right] = \sum_{i \neq i^*} \sum_{j \neq i^*} \mathbb{E}\left[ v_i^T v_j \alpha_i \alpha_j \right]
\end{equation}
By the orthogonality of value representations $\mathbb{E}[v_i^T v_j] = \delta_{ij} \|v_i\|_2^2 = \delta_{ij}$, the cross-terms vanish:
\begin{equation}
\mathbb{E}\left[ \|e\|_2^2 \right] = \sum_{i \neq i^*} \mathbb{E}\left[ \alpha_i^2 \right] = \sum_{i \neq i^*} \left( (\mathbb{E}[\alpha_i])^2 + \operatorname{Var}(\alpha_i) \right)
\end{equation}
Applying Theorem~\ref{thm:transmutation} and Proposition~\ref{thm:fwht}:
\begin{equation}
\mathbb{E}[\alpha_i] = \exp\left( -\frac{1 - \tilde{k}_i^T \tilde{k}_{i^*}}{\tau^2} \right) \le \exp\left( -\frac{1 - \mu_{\max}}{\tau^2} \right)
\end{equation}
and the estimator variance under Orthogonal Random Features is bounded by $\operatorname{Var}(\alpha_i) \le \frac{d_k - 1}{D(d_k + 2)}$. Summing over all $K - 1$ interfering associations:
\begin{equation}
\mathbb{E}\left[ \|e\|_2^2 \right] \le (K - 1) \left[ \exp\left( -\frac{2(1 - \mu_{\max})}{\tau^2} \right) + \frac{d_k - 1}{D(d_k + 2)} \right]
\end{equation}
Applying Jensen's inequality $\mathbb{E}[\mathcal{E}_{\text{copy}}] \le \sqrt{\mathbb{E}[\|e\|_2^2]}$ and using the sub-additivity $\sqrt{a + b} \le \sqrt{a} + \sqrt{b}$ yields Equation~\ref{eq:copy_error_bound}.

\bigskip
\noindent
\textbf{Part (iii): Contrast with Linear DeltaNet.}
In Linear DeltaNet, $D = d_k$ (e.g., $d_k = 32$). When $K > d_k$, Sylvester's rank inequality enforces $\operatorname{rank}([k_1, \dots, k_K]) \le d_k < K$. By the Pigeonhole Principle, at least $K - d_k$ keys must lie in the linear span of previous keys, creating non-trivial co-planar dependencies where $\mu_{\max} \to 1$. Consequently, the exponential suppression term $\exp(-(1 - \mu_{\max})/\tau^2)$ degenerates to $\Omega(1)$, and the variance cannot be suppressed (as $D$ cannot exceed $d_k$). The copying error remains lower-bounded by $\Omega(1)$, destroying induction-head retrieval. In contrast, Kernel-DeltaNet projects into $D = 512 \gg d_k$, ensuring that even for $K = 256$ or $K = 512$, the keys span all $K$ dimensions independently, unlocking near-perfect copying.
\end{proof}
